\section{Gestão de Projeto Ágil}
\label{sec:gestao-projeto-agil}

A organização das atividades de desenvolvimento também exigiu atenção à gestão do projeto,
as prioridades e ao acompanhamento do progresso. Em contextos de software, a gestão ágil busca
reduzir a incerteza por meio de entregas incrementais, feedback frequente e adaptação contínua
às mudanças, em vez de dependência excessiva de planos rígidos e de documentação extensa
\cite{cohn2009}. Nesse sentido, a gestão do módulo foi apoiada por ferramentas e práticas que
permitiram organizar o backlog, acompanhar pendências e controlar o avanço das tarefas de forma
transparente.

\subsection{OpenProject}
\label{subsec:openproject}

O OpenProject é uma plataforma de gestão de projetos e de colaboração que oferece recursos de
backlog, planejamento, acompanhamento de tarefas, cronogramas e relatórios de progresso
\cite{openproject2026}. No estágio, ele foi utilizado para registrar issues, atribuir
responsabilidades, definir prioridades e monitorar o avanço das atividades relacionadas ao
desenvolvimento do módulo. Essa ferramenta contribuiu para manter o trabalho organizado,
reduzindo a perda de contexto e permitindo uma visão clara do status das pendências em cada
etapa do projeto.

\subsection{Scrum}
\label{subsec:scrum}

O Scrum é uma das abordagens ágeis mais conhecidas e organiza o trabalho em ciclos curtos,
chamados \textit{sprints}, com papéis definidos, backlog priorizado e cerimônias de
planejamento, revisão e retrospectiva \cite{cohn2009}. A estrutura do Scrum é especialmente útil
em equipes maiores e em projetos em que a mudança de requisitos é frequente, pois favorece a
entrega incremental de valor e a revisão contínua do produto. No caso do estágio, a metodologia
não foi seguida em sua forma estrita, mas os princípios de priorização, transparência e
acompanhamento incremental foram relevantes para organizar o desenvolvimento do módulo.

\subsection{Extreme Programming (XP)}
\label{subsec:extreme-programming-xp}

O Extreme Programming (XP) enfatiza práticas de engenharia que aumentam a qualidade do software
e a capacidade de responder rapidamente às mudanças de requisitos \cite{beck2004extreme}. Entre
suas práticas mais conhecidas estão pequenas entregas, feedback frequente, desenvolvimento
orientado a testes, integração contínua, refatoração e foco em simplicidade. Como o módulo foi
implementado em um sistema legado, com regras de negócio bem definidas e necessidade de
validação constante, o XP se mostrou mais alinhado ao ritmo de trabalho do que a execução
rígida
de cerimônias do Scrum. As entregas menores, os testes automatizados e a revisão contínua
reduziram o risco de regressão e facilitaram a adaptação à estrutura já existente do Banco de
Especialistas.

A combinação de priorização em backlog, uso do OpenProject para o acompanhamento das atividades
e adoção de práticas do XP permitiu manter um fluxo de trabalho disciplinado, mesmo em um
contexto de desenvolvimento com escopo específico e um único responsável técnico. Essa abordagem
equilibrou organização gerencial e rigor técnico, o que foi fundamental para a construção do
módulo de solicitações.

% TO-DO: Uma tabela comparando Scrum e XP, mostrando as práticas adotadas e as que foram 
% descartadas, com justificativa.
\todo[inline]{TODO:Uma tabela comparando Scrum e XP, mostrando as práticas adotadas e as que foram 
descartadas, com justificativa.}